\documentclass[12pt,a4paper]{article}
\usepackage{amsmath,amsthm,amssymb,amsfonts}
\usepackage[T1]{fontenc}
\usepackage[utf8]{inputenc}
\usepackage[ngerman]{babel}
\usepackage{hyperref}
\usepackage{geometry}
\geometry{margin=2.5cm}

\newtheorem{theorem}{Satz}[section]
\newtheorem{lemma}[theorem]{Lemma}
\newtheorem{proposition}[theorem]{Proposition}
\newtheorem{corollary}[theorem]{Korollar}
\newtheorem{definition}[theorem]{Definition}
\newtheorem{remark}[theorem]{Bemerkung}
\newtheorem{conjecture}[theorem]{Vermutung}

\begin{document}

\title{Taos probabilistischer Ansatz zur Collatz-Vermutung:\\
Fast alle Bahnen erreichen fast beschränkte Werte}
\author{Michael Fuhrmann}
\date{März 2026}
\maketitle

\begin{abstract}
Im Jahr 2022 veröffentlichte Terence Tao ein bahnbrechendes Ergebnis zur Collatz-Vermutung:
Für jede Funktion $f\colon \mathbb{N} \to \mathbb{R}_{>0}$ mit $f(n) \to \infty$ haben
fast alle positiven ganzen Zahlen $n$ (im Sinne der logarithmischen Dichte) schließlich
ein Collatz-Iteriertes unterhalb von $f(n)$. Diese Arbeit präsentiert Taos Ansatz im
Detail: das probabilistische Zufallswalk-Modell für die Syrakus-Funktion, die
Fourier-analytischen Hilfsmittel zur Behandlung von Korrelationen zwischen Schritten,
den Rahmen der logarithmischen Dichte und die Schlüsselabschätzung, die zeigt, dass
divergente Collatz-Bahnen die logarithmische Dichte null haben. Wir diskutieren auch,
warum Taos Methode nicht ausreicht, um die vollständige Vermutung zu beweisen, und
was nötig wäre, um die verbleibende Lücke zu schließen.
\end{abstract}

\tableofcontents

\newpage

\section{Einleitung}

Terence Taos Arbeit von 2022 ``Almost all Collatz orbits attain almost bounded values''
\cite{Tao2022}, veröffentlicht in \emph{Forum of Mathematics, Sigma}, stellt den bedeutendsten
Fortschritt bei der Collatz-Vermutung seit Jahrzehnten dar. Sie beweist die Vermutung
nicht, liefert aber das stärkste bekannte unbedingte Resultat zu ihren Gunsten.

Die entscheidende Erkenntnis ist probabilistischer Natur: Die Syrakus-Funktion
$S(n) = (3n+1)/2^{\nu_2(3n+1)}$ verhält sich wie ein ``zufälliger'' Spaziergang mit
starkem Abwärtsdrift, und Tao quantifiziert diesen Drift präzise mithilfe der
Fourier-Analysis über die 2-adischen ganzen Zahlen.

\subsection{Hauptresultat}

\begin{theorem}[Tao 2022]\label{thm:tao_hauptsatz}
Für jede Funktion $f\colon \mathbb{N} \to \mathbb{R}_{>0}$ mit $f(n) \to \infty$ für
$n \to \infty$ gilt:
\[
  \mathrm{logDichte}\bigl(\{ n \in \mathbb{N} : \exists k \geq 0,\; T^k(n) \leq f(n) \}\bigr) = 1.
\]
Äquivalent: Für jedes $\varepsilon > 0$ und jedes $f \to \infty$ hat die Menge der $n$,
für die kein Iteriertes von $T$ unter $f(n)$ fällt, die logarithmische Dichte 0.
\end{theorem}

\section{Das Syrakus-Zufallswalk-Modell}

\begin{definition}[Logarithmische Dichte]\label{def:logdichte}
Die \emph{logarithmische Dichte} einer Menge $A \subseteq \mathbb{N}$ ist
\[
  \mathrm{logDichte}(A) \;=\; \lim_{N \to \infty} \frac{1}{\log N} \sum_{\substack{n=1 \\ n \in A}}^{N} \frac{1}{n},
\]
sofern der Grenzwert existiert. Die logarithmische Dichte wird durch die natürliche Dichte
dominiert, ist jedoch empfindlicher für die Verteilung kleiner Primzahlen.
\end{definition}

\begin{definition}[Syrakus-Funktion]\label{def:syrakus}
Für eine ungerade positive ganze Zahl $m$ ist die \emph{Syrakus-Funktion}
\[
  S(m) \;=\; \frac{3m+1}{2^{\nu_2(3m+1)}},
\]
wobei $\nu_2(3m+1)$ die 2-adische Bewertung bezeichnet (die Anzahl, wie oft 2 in $3m+1$ aufgeht).
Der Wert $E(m) := \nu_2(3m+1)$ ist die Anzahl der ``geraden Schritte'', die beim
Übergang von $m$ zu $S(m)$ in der Standard-Collatz-Iteration verbraucht werden.
\end{definition}

\begin{proposition}[Verteilung von $E(m)$]\label{prop:vert_E}
Für eine zufällige ungerade positive ganze Zahl $m$ folgt die Zufallsvariable
$E(m) = \nu_2(3m+1)$ näherungsweise einer geometrischen Verteilung mit Parameter $1/2$:
\[
  \Pr[E(m) = k] \;\approx\; 2^{-k} \quad \text{für } k \geq 1.
\]
Genauer: Unter den ungeraden ganzen Zahlen $m \in [1, N]$ beträgt der Anteil mit
$E(m) = k$ den Wert $2^{-k}(1 + O(N^{-1/2}))$ für $N \to \infty$.
\end{proposition}

\begin{proof}
Für eine ungerade ganze Zahl $m$ gilt $m \equiv 1$ oder $3 \pmod{4}$.
Falls $m \equiv 3 \pmod{4}$: $3m+1 \equiv 10 \equiv 2 \pmod{4}$, also $\nu_2(3m+1) = 1$, d.h.\ $E(m)=1$.
Falls $m \equiv 1 \pmod{4}$: $3m+1 \equiv 4 \equiv 0 \pmod{4}$, also $\nu_2(3m+1) \geq 2$, d.h.\ $E(m)\geq 2$.
Also gilt $E(m) = 1$ genau dann, wenn $m \equiv 3 \pmod{4}$: Wahrscheinlichkeit $1/2$ unter ungeraden Zahlen.
Allgemein: $E(m) \geq k$ genau dann, wenn $2^k \mid 3m+1$, was (da $\gcd(3,2)=1$) äquivalent ist zu
$m \equiv r_k \pmod{2^k}$ für einen eindeutig bestimmten ungeraden Rest $r_k$.
Dies tritt mit Wahrscheinlichkeit $2^{-(k-1)}$ unter ungeraden Zahlen auf.
Damit: $\Pr[E(m) = k] = \Pr[E(m) \geq k] - \Pr[E(m) \geq k+1] = 2^{-(k-1)} - 2^{-k} = 2^{-k}$.
\end{proof}

\begin{definition}[Logarithmischer Zufallswalk]\label{def:zufallswalk}
Um die ``Größe'' der Syrakus-Iterationen zu modellieren, wird der natürliche Logarithmus
($\log = \ln$) genommen. Für ungerades $m$ gilt:
\[
  \ln S(m) \;=\; \ln m + \ln 3 - E(m) \ln 2 + O(1/m).
\]
Die Inkremente $X_k = \ln 3 - E_k \ln 2$ (wobei $E_k = E(S^{k-1}(m))$) bilden die
einzelnen Schritte des Zufallswalks. Der Erwartungswert ist:
\[
  \mathbb{E}[X_k] \;=\; \ln 3 - \ln 2 \cdot \mathbb{E}[E_k] = \ln 3 - 2\ln 2 = \ln(3/4) < 0.
\]
Da $\ln(3/4) \approx -0{,}288 < 0$, hat der Walk einen strikten negativen Drift.
In der Basis-2-Schreibweise gilt äquivalent $\log_2(3/4) = \log_2 3 - 2 \approx -0{,}415$
pro Syrakus-Schritt; beide Formulierungen beschreiben denselben Abwärtstrend.
\end{definition}

\begin{theorem}[Negativer Drift impliziert fast sichere Konvergenz]\label{thm:drift}
Seien $(X_k)_{k \geq 1}$ i.i.d.-Zufallsvariablen mit $\mathbb{E}[X_1] = \ln(3/4) < 0$.
Dann gilt für den Zufallswalk $W_n = \sum_{k=1}^n X_k$: $W_n \to -\infty$ fast sicher.
Insbesondere fällt der Walk für jede feste Schranke $L$ mit Wahrscheinlichkeit 1
schließlich unter $L$.
\end{theorem}

\begin{proof}
Dies folgt unmittelbar aus dem starken Gesetz der großen Zahlen: $W_n/n \to \mathbb{E}[X_1] = \ln(3/4) < 0$
fast sicher. Da $\mathbb{E}[X_1] < 0$, gilt $W_n \to -\infty$ f.s.
\end{proof}

\begin{remark}\label{rem:nicht_unabhaengig}
Die entscheidende Schwierigkeit liegt darin, dass die Schritte $X_k = \ln 3 - E(S^{k-1}(m)) \ln 2$
\emph{nicht} unabhängig sind: die Verteilung von $E(S^k(m))$ hängt vom Wert von $S^k(m)$
ab, der wiederum von der gesamten Bahn-Geschichte abhängt. Taos Hauptbeitrag besteht
darin, diese Abhängigkeiten mittels Fourier-Analysis zu behandeln.
\end{remark}

\section{Fourier-Analysis über den 2-adischen ganzen Zahlen}

\begin{definition}[Charaktere von $\mathbb{Z}/2^k\mathbb{Z}$]\label{def:charaktere}
Für $k \geq 1$ und $\xi \in \mathbb{Z}/2^k\mathbb{Z}$ ist der \emph{Charakter}
$\chi_{\xi,k}\colon \mathbb{Z}/2^k\mathbb{Z} \to \mathbb{C}^*$ definiert durch
\[
  \chi_{\xi,k}(n) \;=\; e^{2\pi i \xi n / 2^k}.
\]
Die 2-adische Fourier-Transformierte einer Funktion $f\colon \mathbb{Z}_2 \to \mathbb{C}$
auf Resten mod $2^k$ ist $\hat{f}(\xi) = \sum_{n \bmod 2^k} f(n) \chi_{\xi,k}(n)$.
\end{definition}

\begin{definition}[Syrakus-Exponentialsumme]\label{def:exp_summe}
Für $\xi \in \mathbb{R}/\mathbb{Z}$ und $k \geq 1$ ist die \emph{Syrakus-Exponentialsumme}
\[
  F_k(\xi) \;=\; \mathbb{E}_{m \bmod 2^k,\, m \text{ ungerade}}\bigl[ e^{2\pi i \xi \log_2 S(m)} \bigr],
\]
wobei der Erwartungswert über ungerade Reste mod $2^k$ gebildet wird.
\end{definition}

\begin{lemma}[Taos Schlüsselabschätzung]\label{lem:schluessel}
Für alle $\xi \in \mathbb{R}/\mathbb{Z}$ mit $\|\xi\|_{\mathbb{R}/\mathbb{Z}} \geq \delta > 0$
existiert $c(\delta) > 0$ mit
\[
  |F_k(\xi)| \;\leq\; e^{-c(\delta) k}
\]
für alle hinreichend großen $k$. Dieser exponentielle Abfall in $k$ ist der Kern
von Taos Beweis.
\end{lemma}

\begin{proof}[Beweisansatz (informell -- kein rigoroser Beweis)]
% BUG-B7-P31-EN/DE-001 behoben: Klar als informal deklariert; ξ-Bedingung korrigiert
\emph{Hinweis: Der folgende Text ist nur eine heuristische Skizze der i.i.d.-Näherung.
Ein rigoroser Beweis erfordert Taos Entkopplungsargument; siehe~\cite{Tao2022}.
Der Übergang von Summe zu Produkt ist nicht unmittelbar, und die
Unabhängigkeitsannahme ist nicht exakt erfüllt.}

Die Exponentialsumme $F_k(\xi)$ faktorisiert \emph{näherungsweise} als Produkt über
``unabhängige'' Schritte. Wären die $E_j$ exakt i.i.d.\ geometrisch(1/2) verteilt, so ergäbe sich:
\[
  F_k(\xi) = \left(\sum_{j=1}^\infty 2^{-j} e^{2\pi i \xi (\log_2 3 - j)}\right)^k,
\]
und die innere Summe ist eine geometrische Reihe in $e^{-2\pi i \xi}$:
\[
  F_k(\xi) = \left(\frac{e^{2\pi i \xi \log_2 3}}{2 - e^{-2\pi i \xi}}\right)^k.
\]
Für $\|\xi\|_{\mathbb{R}/\mathbb{Z}} > 0$ (d.\,h.\ $\xi \notin \mathbb{Z}$, was durch
die Lemma-Voraussetzung $\|\xi\|_{\mathbb{R}/\mathbb{Z}} \geq \delta > 0$ garantiert ist),
hat der Nenner $2 - e^{-2\pi i \xi}$ Betrag $< 2$, also ist der Betrag des inneren
Faktors $< 1$, was exponentiellen Abfall in $k$ liefert.
Tao behandelt die Nicht-Unabhängigkeit mittels eines Entkopplungsarguments, das die
Fourier-Struktur von $\mathbb{Z}_2$ nutzt und dem informellen Modell Rigorosität verleiht.
\end{proof}

\section{Das Argument der logarithmischen Dichte}

\begin{definition}[Divergente Menge]\label{def:divergent}
Definiere die \emph{divergente Menge} als
\[
  \mathcal{D} \;=\; \bigl\{ n \in \mathbb{N} : T^k(n) \geq f(n) \text{ für alle } k \geq 0 \bigr\},
\]
d.h.\ die Menge der ganzen Zahlen, deren Collatz-Bahn nie unter $f(n)$ fällt.
Taos Satz behauptet $\mathrm{logDichte}(\mathcal{D}) = 0$.
\end{definition}

\begin{theorem}[Logarithmische Dichte divergenter Bahnen]\label{thm:logdichte_null}
Für jedes $f\colon \mathbb{N} \to \mathbb{R}_{>0}$ mit $f(n) \to \infty$ hat die
Menge $\mathcal{D}$ die logarithmische Dichte null.
\end{theorem}

\begin{proof}[Beweisüberblick (nach Tao \cite{Tao2022})]
Der Beweis verläuft in drei Hauptschritten.

\textbf{Schritt 1: Reduktion auf die Syrakus-Abbildung.}
Man arbeitet mit der komprimierten Syrakus-Funktion $S$ auf ungeraden Zahlen. Das
Collatz-Iterierte $T^k(n)$ ist höchstens $f(n)$ genau dann, wenn $S^j(n) \leq C \cdot f(n)$
für ein $j$ und eine Konstante $C$ gilt. Es genügt also zu zeigen, dass die logarithmische
Dichte der ungeraden $n$, für die alle Syrakus-Iterationen über $f(n)$ bleiben, null ist.

\textbf{Schritt 2: Diskretisierung via 2-adische Reste.}
Für jede Skala $N$ untersucht man die Bahn von $n$ modulo $2^k$ für $k = O(\log \log N)$.
Die Syrakus-Funktion wirkt auf Resten mod $2^k$, und die Fourier-Abschätzungen aus
Lemma~\ref{lem:schluessel} zeigen, dass Korrelationen zwischen Bahnschritten exponentiell
in $k$ abnehmen.

\textbf{Schritt 3: Exponentialsummen-Abschätzung.}
Man schreibt den Indikator für ``Bahn bleibt über $f(n)$'' als Fourier-Integral, wendet
Lemma~\ref{lem:schluessel} an, um die Fourier-Koeffizienten zu beschränken, und integriert.
Die Schranke $|F_k(\xi)| \leq e^{-c(\delta)k}$ impliziert, dass die logarithmische Dichte
der ganzen Zahlen mit nicht-absteigenden Bahnen (nach $k$ Syrakus-Schritten) höchstens
$C e^{-c(\delta)k}$ beträgt, was für $k \to \infty$ gegen 0 strebt.
\end{proof}

\begin{corollary}[Collatz-Bahnen sind fast beschränkt]\label{cor:fast_beschraenkt}
Für jedes $\varepsilon > 0$ hat die Menge der $n \in \mathbb{N}$, für die $T^k(n) > n^\varepsilon$
für alle $k \geq 0$ gilt, die logarithmische Dichte null. Insbesondere enthalten fast alle
(im Sinne der logarithmischen Dichte) Collatz-Bahnen beliebig kleine Werte relativ zum
Ausgangspunkt.
\end{corollary}

\section{Die Rolle der Ergodentheorie}

\begin{definition}[Syrakus-Zufallswalk als ergodischer Prozess]\label{def:ergodisch_rw}
Die Folge der Log-Verhältnisse $(\log S^k(n) / \log n)_{k \geq 0}$ kann als Trajektorie
eines Punktes in einem maßerhaltenden dynamischen System betrachtet werden.
\textbf{Wichtiger Hinweis:} Die Syrakus-Abbildung $S$ ist \emph{nicht} maßerhaltend
bezüglich des Haar-Maßes auf $\mathbb{Z}_2^{\text{ungerade}}$.
Um den Birkhoff-Ergodensatz anwenden zu können, muss ein separat konstruiertes
$S$-invariantes Maß $\nu$ auf $\mathbb{Z}_2^{\text{ungerade}}$ verwendet werden
(wie in \cite{Tao2022} und \cite{Kontorovich2010} rigoros durchgeführt).
Unter einem solchen invarianten Maß $\nu$ ist die Funktion
$g(x) = \log_2 3 - E(x)$ (wobei $E(x) = \nu_2(3x+1)$) die ``Schrittweiten''-Funktion
mit $\nu$-Erwartungswert
\[
  \int g\, d\nu = \log_2 3 - \mathbb{E}_\nu[E].
\]
Wird $\nu$ so gewählt, dass $E$ eine geometrische Verteilung mit Parameter $1/2$ besitzt
(was im i.i.d.-Näherungsmodell gilt), so ist $\mathbb{E}_\nu[E] = 2$ und
\[
  \frac{1}{k} \sum_{j=0}^{k-1} g(S^j(x)) \;\longrightarrow\; \log_2 3 - 2 = \log_2(3/4) < 0
\]
für $\nu$-fast alle $x$ nach dem Birkhoff-Ergodensatz.
\end{definition}

\begin{theorem}[Ergodische Konvergenz von Syrakus-Bahnen]\label{thm:ergod_konv}
Sei $\nu$ ein $S$-invariantes ergodisches Maß auf $\mathbb{Z}_2^{\text{ungerade}}$, unter
dem $E(x) = \nu_2(3x+1)$ den Erwartungswert 2 besitzt (wie im i.i.d.-geometrisch$(1/2)$-Modell).
Dann gilt für $\nu$-fast alle $x$:
\[
  \lim_{k \to \infty} \frac{\log S^k(x)}{k} \;=\; \log(3/4) < 0.
\]
Dies bedeutet, dass der Logarithmus der Bahn linear abnimmt, d.h.\ die Bahn fällt
\emph{geometrisch schnell} im absoluten Sinne.

\textbf{Hinweis:} Das Haar-Maß auf $\mathbb{Z}_2^{\text{ungerade}}$ ist \emph{nicht}
$S$-invariant, weshalb der Birkhoff-Ergodensatz nicht direkt mit dem Haar-Maß angewendet
werden darf. Die obige Aussage verwendet das invariante Maß $\nu$; eine explizite
Konstruktion findet sich in \cite{Kontorovich2010}.
\end{theorem}

\begin{proof}
Das Integral $\int g\, d\nu = \mathbb{E}_\nu[\log_2 3 - E] = \log_2 3 - \mathbb{E}_\nu[E]$.
Nach Voraussetzung gilt $\mathbb{E}_\nu[E] = 2$
(aus Proposition~\ref{prop:vert_E} im i.i.d.-Modell).
Also ist der $\nu$-ergodische Mittelwert von $g$ gleich $\log_2 3 - 2 = \log_2(3/4) \approx -0{,}415 < 0$.
Da $\nu$ $S$-invariant und ergodisch ist, liefert der Birkhoff-Ergodensatz die behauptete
Konvergenz für $\nu$-fast alle $x$.
\end{proof}

\begin{lemma}[Korrespondenzprinzip im Collatz-Kontext]\label{lem:furstenberg}
Sei $(X, \mathcal{B}, \mu, T)$ ein maßerhaltendes dynamisches System und $A \in \mathcal{B}$
mit $\mu(A) > 0$. Ein Korrespondenzprinzip (vgl.~\cite{Tao2022}) stellt eine
strukturelle Verbindung zwischen kombinatorischen Eigenschaften der Collatz-Bahn und der
ergodischen Struktur von $(\mathbb{Z}_2, S, \mu)$ her.
Im Collatz-Kontext dient das Prinzip dazu, Dichtaussagen aus dem 2-adischen Rahmen auf
Aussagen über die natürliche Dichte in $\mathbb{N}$ zu übertragen.
\end{lemma}

\section{Warum Taos Resultat die Vermutung nicht beweist}

\begin{remark}[Die Lücke zwischen logarithmischer Dichte und vollständiger Vermutung]\label{rem:luecke}
Taos Satz zeigt, dass die Menge der ``schlecht verhaltenen'' Bahnen die logarithmische
Dichte 0 hat. Die vollständige Collatz-Vermutung erfordert jedoch zu zeigen, dass diese
Menge \emph{leer} ist (oder äquivalent, dass jede Bahn schließlich 1 erreicht). Das sind
sehr unterschiedliche Aussagen: Eine Menge der logarithmischen Dichte 0 kann trotzdem
unendlich sein (in der Tat: natürliche Dichte 0 aber unendlich groß).
\end{remark}

\begin{proposition}[Hindernisse für einen vollständigen Beweis]\label{prop:hindernisse}
Folgende Hindernisse verhindern, dass Taos Ansatz die vollständige Vermutung beweist:
\begin{enumerate}
  \item \textbf{Individuell vs.\ statistisch}: Tao kontrolliert das durchschnittliche
        Verhalten, aber nicht das Verhalten einer einzelnen Bahn.
  \item \textbf{Keine Rekurrenz}: Die 2-adische Fourier-Methode liefert Dichte-Null-Resultate,
        aber keine Rekurrenz für jeden Startpunkt.
  \item \textbf{Ausschluss von Zyklen}: Taos Argument schließt nicht die Existenz
        extrem großer Zyklen aus, die manche Bahnen ``einfangen'' könnten.
  \item \textbf{Die $+1$-Barriere}: Das $+1$ in $3n+1$ erzeugt arithmetische Korrelationen,
        die vollständige Unabhängigkeit der Bahnschritte verhindern und die Stärke des
        Fourier-Ansatzes begrenzen.
\end{enumerate}
\end{proposition}

\begin{conjecture}[Collatz-Vermutung, vollständige Form]\label{conj:vollstaendig}
Für jedes $n \in \mathbb{N}$ existiert $k \geq 0$ mit $T^k(n) = 1$.
\end{conjecture}

\section{Diskussion und offene Probleme}

Taos Resultat ist in mehrfacher Hinsicht bedeutsam:
\begin{enumerate}
  \item Es stellt die stärkste bekannte quantitative Aussage zur Vermutung auf.
  \item Es zeigt, dass Gegenbeispiele --- falls sie existieren --- eine sehr ``dünne''
        Menge bilden müssen (logarithmische Dichte 0).
  \item Es stellt einen Werkzeugkasten (Fourier-Analysis über $\mathbb{Z}_2$,
        Syrakus-Zufallswalks) bereit, der weiter entwickelt werden kann.
\end{enumerate}

Offene Probleme:
\begin{enumerate}
  \item Kann das logarithmische Dichte-Resultat auf natürliche Dichte verbessert werden?
  \item Kann man zeigen, dass divergente Bahnen natürliche Dichte 0 haben?
  \item Kann der Fourier-analytische Ansatz mit $p$-adischen algebraischen Methoden
        kombiniert werden, um alle nicht-trivialen Zyklen auszuschließen?
  \item Gibt es einen ``strukturellen'' Grund (jenseits probabilistischer Heuristiken),
        warum das $+1$ alle Bahnen zur Konvergenz zwingt?
\end{enumerate}

\begin{thebibliography}{99}

\bibitem{Tao2022}
T.~Tao.
\newblock Almost all Collatz orbits attain almost bounded values.
\newblock \emph{Forum of Mathematics, Sigma}, 10:e72, 2022.

\bibitem{Collatz1937}
L.~Collatz.
\newblock Unveröffentlichtes Problem, vorgestellt beim Internationalen Mathematikerkongress, 1937.

\bibitem{Lagarias1985}
J.~C. Lagarias.
\newblock The $3x+1$ problem and its generalizations.
\newblock \emph{American Mathematical Monthly}, 92(1):3--23, 1985.

\bibitem{Lagarias2010}
J.~C. Lagarias (Hrsg.).
\newblock \emph{The Ultimate Challenge: The $3x+1$ Problem}.
\newblock American Mathematical Society, Providence, RI, 2010.

\bibitem{Terras1976}
R.~Terras.
\newblock A stopping time problem on the positive integers.
\newblock \emph{Acta Arithmetica}, 30(3):241--252, 1976.

\bibitem{Koblitz1984}
N.~Koblitz.
\newblock \emph{$p$-adic Numbers, $p$-adic Analysis, and Zeta Functions}.
\newblock Springer-Verlag, New York, 1984.

\bibitem{Kontorovich2010}
A.~Kontorovich und J.~C. Lagarias.
\newblock Stochastic models for the $3x+1$ and $5x+1$ problems and related problems.
\newblock In \emph{The Ultimate Challenge}, S.\ 131--188, AMS, 2010.

\end{thebibliography}

\end{document}
